% Ad Familiares 2.12 --- headnote
% ad-familiares-02-12, just before 26 June 50 BC, Cilician camp

\textit{Cicero to M. Caelius Rufus, curule aedile,
written from the Cilician camp shortly before the sixth
day before the Kalends of Quintilis (26 June) 50 BC
(Perseus dateline: \textit{Scr. in castris in Cilicia
paulo ante vi K. Quint. a. 704 (50)}). The setting is
the summer encampment Cicero had mentioned in
\textit{Fam} 2.13 as the last station before his
intended departure: news from Rome has been slow and
fragmentary, and the disorder of the spring's public
meetings is still reaching him by old report rather
than by Caelius's customary newsletter.}

\textit{The middle section is the most famous of all
the homesick passages in the letters. Cicero has just
parted at Pessinus from two friends of Caelius's,
Diogenes and Philo, who are pressing on into the
Anatolian interior to Adiatorix --- and the contrast
between their road and his own breaks open into the
\textit{Urbem, mi Rufe} peroration: hold to the City,
live in her light, all foreign service is squalid and
obscure for men whose industry can shine at Rome.
The closing section converts the same feeling into
politics: he hopes he has won the praise of integrity,
and the modest possibility of a triumph is not worth
the months of separation it would cost. A note in three
movements --- anxious, longing, weary --- closing with
the standing request for Caelius's newsletter to meet
him on the homeward road.}
